% !TeX root = ../tesis.tex
%%%%%%%%%%%%%%%%%%%%%%%%%%%%%%%%%%%%%%%%%%%%%%%%%%%%%%%%%%%%%%%%%%%
% SECCIÓN 7.1 — EJE TECNOLÓGICO (Issue #23)
% Papel en el hilo: el instrumento que hace medible la estructura.
% Esquema: problema (1985) → qué se hizo → replicabilidad → aporte →
% límites y transición a 7.2. Plan: CONTEXTO_CAP7.md §4.
% Fuera a propósito: T_b fijo (#32), grafo simple (D9), nodo
% «Memorial», 13 nodos del snapping; TransCAD/NAVTEQ sin verificar.
%%%%%%%%%%%%%%%%%%%%%%%%%%%%%%%%%%%%%%%%%%%%%%%%%%%%%%%%%%%%%%%%%%%
\section{Eje Tecnológico: un Instrumento para Medir la Estructura}
\label{sec:conclusiones-tecnologico}

El Programa Maestro del Metro de 1985 descartó los dos corredores del
Periférico porque no alcanzaban la captación que pedía su red
(§\ref{subsubsec:anillo-pm-2010}). Sus parámetros salían de la demanda, y con
ellos la evaluación no tenía forma de medir lo que el anillo aportaba a la
estructura de la red. Aquel estudio se hizo con otras tecnologías, las de su
momento. Medir la estructura pide hoy otro tipo de herramienta, y por eso el
primer resultado de esta tesis es un instrumento: una cadena abierta que mide
la estructura de la red de transporte con datos públicos.

La cadena tiene tres capas, y cada una cumple una tarea
(Figura~\ref{fig:macro-arquitectura-stack}). \texttt{Apímetro} reúne la
información de doce sistemas de transporte y 141 demarcaciones, que cada
organismo publica a su manera, y la entrega en un solo formato. Con ello cubre
algo que la Tarjeta MI dejó pendiente en 2019, cuando integró la tarifa pero no
los datos (§\ref{subsubsec:espectro-cierre}). \texttt{VFTModel} arma con esa
información un grafo de 11,115 nodos y 45,517 aristas, asigna a cada arista un
tiempo con la ecuación VFT y calcula los ocho indicadores
(§\ref{sec:vft-indicadores}). \texttt{Transport-gis-zmvm-mjg}, por último,
lleva los resultados a mapas de QGIS y a \termino{dashboards} de Tableau, que
es donde un urbanista puede leerlos.

Que la cadena sirva como instrumento, y no solo para resolver un caso, se
comprobó en el Capítulo~\ref{ch:analisis-red-anillo}. Los tres escenarios se
procesaron con el mismo código. Para sumar el anillo bastó registrarlo en
\texttt{Apímetro} como un sistema más, y cada caso corrió en sus propios
contenedores, sin que los datos de uno alteraran los de otro
(§\ref{subsec:anillo-procesamiento}). Con esa separación se pueden comparar
corredores y estimar el efecto de una línea nueva antes de construirla
(§\ref{sec:vft-salidas}). Es el mismo ejercicio que hizo el Programa Maestro al
simular su red varias veces, solo que ahora lo que se compara es la estructura.

El aporte de este eje no se limita a la \zmvm{}. La cadena no depende de
licencias comerciales y se puede repetir en cualquier ciudad que tenga
\termino{feeds} \gtfs{} y cartografía de base, algo que ya cumple la mayoría de las ciudades
latinoamericanas de tamaño parecido \citep{barrington2017world}. El código se
publica bajo la \textit{Business Source License~1.1}, que pasa a
\textit{Apache~2.0} en abril de 2029, y viene con una wiki, un foro,
\texttt{seed.sql} y Docker Compose para que lo use también quien no programa
(§\ref{subsec:futuro-apimetro}). \citet{gurstein2011open} advierte que los
datos abiertos solo empoderan a quien puede procesarlos; una herramienta así
acorta esa distancia. En ese sentido entiende esta tesis la soberanía
tecnológica: una capacidad de análisis que se queda en la institución pública y
no se compra a un proveedor.

El instrumento tiene límites, y conviene tenerlos presentes al leer sus
resultados. Las velocidades de superficie vienen de \texttt{Apímetro} y no se
han contrastado con velocidades observadas por ruta
(§\ref{subsec:validacion-impedancia}). La caída de eficiencia $\Delta E$ se
calculó solo como ejercicio exploratorio, con una corrida anterior a la vigente
(§\ref{subsec:vulnerabilidad-estructural}). Las cinco estaciones del Tren
Interurbano quedan fuera de la componente gigante. El límite principal es otro:
el modelo mide la estructura, no el uso. No calcula demanda ni capacidad, aunque
\texttt{Apímetro} ya guarda la capacidad de los vehículos. Estos pendientes se
retoman como trabajo futuro en §\ref{sec:conclusiones-academico}. Con el
instrumento probado, falta ver qué dice esa medición sobre la red, que es el
tema de §\ref{sec:conclusiones-urbano}.
